\section{Interpretation and Business Insight}
\label{sec:insight}

\subsection{What the models rely on}

Importance is measured by permutation on the held-out window rather than by
impurity on the training set: impurity-based scores are computed where the model
was fitted and inflate high-cardinality features \citep{strobl2007importance},
which with 74 product categories in play is not hypothetical
(Figure~\ref{fig:importance}).

\textbf{Problem 2 has one dominant driver, and it is actionable.} How late the
parcel was carries 0.206 of PR-AUC on its own, how long it took a further 0.074,
the size of the order 0.045; nothing else clears 0.011. The top two account for
\textbf{76.9\%} of total positive importance. That is the honest description ---
the model asks how late the parcel was and how big the order was --- and it
explains why a single depth-limited tree captured 96\% of the achievable signal.
For the customer-experience team that is good news: the driver is something the
business already measures and can already act on.

\textbf{Problem 1A's importances must be read with care.} No single feature
carries much --- the largest, \code{lead\_payment\_type}, costs 0.034 days of
MAE, and the top two together only 30\% of total positive importance --- and that
is a property of the measurement. Permutation importance systematically
understates correlated predictors, because permuting one leaves its substitutes
intact, and \code{estimated\_days}, \code{hist\_route\_lead\_days},
\code{hist\_seller\_lead\_days} and \code{log\_distance} all encode overlapping
information about how far and how long; the selected model is additionally a
blend, so permuting a column degrades only part of it. The reading that survives
is negative and worth having: this model has no single point of failure.

\subsection{From a probability to a decision --- and when not to act}
\label{sec:cost}

A probability is not a decision. Turning one into the other needs a cost, and
Olist publishes none --- so what follows states its assumptions openly and treats
the conclusions as shapes rather than numbers \citep{elkan2001cost}. The
arithmetic is one line: acting is worth it only where
$\text{precision} > c_{\text{action}} / c_{\text{miss}}$.

\begin{table}[htbp]
\centering
\caption{The decision layer under stated (not measured) costs. Break-even
precision is what decides whether a flag is worth acting on at all.}
\label{tab:cost}
\small
\begin{tabular}{@{}lrrrl@{}}
\toprule
\textbf{Flag} & \textbf{$c_{\text{action}}$ : $c_{\text{miss}}$} &
\textbf{break-even} & \textbf{best achieved} & \textbf{verdict} \\
\midrule
Low-review risk, at delivery & R\$5 : R\$40  & 0.125 & 0.37 at the optimum & \best{act} \\
Late-order flag, at checkout & R\$12 : R\$40 & 0.300 & 0.29 at any volume & do not expedite \\
\bottomrule
\end{tabular}
\end{table}

\textbf{The review flag pays, comfortably.} Contacting a customer who was never
going to complain costs perhaps R\$5 of agent time; a 1--2 star review that could
have been headed off costs perhaps R\$40 --- standing in for the damage
\citet{deshpande2023logistics} trace through to seller sales. Break-even
precision is 0.125 and the model reaches 0.369 at its expected-benefit optimum,
flagging 1{,}927 of 12{,}855 delivered orders (15.0\%) of which 712 do go on to
leave a low review, for a net \textbf{R\$1{,}466 per 1{,}000 delivered orders}
(Figure~\ref{fig:threshold}). The margin is wide enough that the conclusion
survives a wrong cost estimate. Capacity changes how the optimum is used rather
than where it is: a team that can work only 5\% of deliveries still captures
R\$1{,}107 of those R\$1{,}466 per thousand orders (Table~\ref{tab:capacity}),
because the top of the list is where the precision is, so benefit per contact
is highest at the smallest capacities and the list is simply worked top-down
until the agents run out.

\textbf{The late-order flag does not pay --- for an expensive action.} This is the
more useful finding, and reporting discrimination alone would have missed it.
Expediting freight costs perhaps R\$12 against the same R\$40 of downstream
damage, so break-even precision is 0.300 --- and the derived flag never reaches
it: precision plateaus at \textbf{0.29} and falls as volume rises, so expected net
benefit is negative at every threshold flagging a meaningful number of orders. A
model with ROC-AUC 0.782 and 2.9$\times$ the no-skill floor is genuinely
informative and still not worth paying R\$12 an order to act on. That changes
\emph{which} action it justifies rather than making it useless: break-even
falls with the cost of the intervention, and two of the available interventions
-- re-quoting the date and pre-warning the buyer -- cost almost nothing
(Section~\ref{sec:stakeholders}).

\subsection{What each stakeholder does with this}
\label{sec:stakeholders}

\textbf{Operations, at checkout} works the four levers of Table~\ref{tab:levers}
in order of cost. Every order gets the re-quoted date: the model's estimate is
off by 3.90 days against the promise's 9.79, so the date shown at checkout can
be tightened materially, which \citet{cui2024promise} show trades directly
against returns and retention. Every order above the recall-oriented threshold
of 0.11 gets a proactive notice, which covers five in six of those that will
actually miss. Expedited freight is reserved for a manual review of the top one
or two per cent of the late list, where precision is 0.27 -- informative, but
under the break-even an automatic rule would need -- and that policy is
revisited only if expediting falls below about R\$11 an order. And the
track-record features and the weekly late series are the monitoring
instruments: a seller whose \code{hist\_seller\_late\_rate} climbs, or a week in
which the carrier's transit time lengthens while the promise does not, is the
February--March 2018 pattern beginning again (Appendix~\ref{app:lateness}).
\textbf{Customer experience, at delivery} receives each day's deliveries ranked
by $p_{\text{low}}$ and works the list from the top: at ten contacts a day --
4.8\% of deliveries -- 57\% of contacts reach a customer who would have left a
1--2 star review and a quarter of all such reviews are reached before they are
written; a team with twice the capacity reaches 38\% of them at 44\% precision
(Table~\ref{tab:capacity}). The dashboard's batch scorer returns exactly this
list, ranked, with the top-$K$ marked for whatever capacity is entered.

These are not two systems: the order flagged at checkout is the order that
arrives late, and the order that arrives late is the one flagged for outreach, so
the first gives operations a chance to prevent the cost the second is left
measuring.
